\begin{tikzpicture}[scale=2]


% === Parâmetros ===
\def\r{3}   % raio do círculo
\def\a{2}   % raio do cilindro

% === Cálculos automáticos ===
\pgfmathsetmacro{\yint}{sqrt(\r*\r - \a*\a)} % coordenada y de interseção
\pgfmathtruncatemacro{\yintb}{\r*\r - \a*\a} % O que vai dentro da raiz, pra impressão
\pgfmathsetmacro{\ang}{atan(\yint/\a)}       % ângulo correspondente


  \begin{axis}[      
      xlabel={$y$}, ylabel={$z$},   
      axis equal image,             % mantém escala x=y para círculos não ficarem ovalados
      every axis x label/.style={at={(ticklabel* cs:0.97)}, anchor=north},
      every axis y label/.style={at={(ticklabel* cs:0.97)}, anchor=east},
      xmin=-1.2*\r, xmax=1.2*\r, ymin=-1.2*\r, ymax=1.2*\r,
      % width=8cm, height=6cm,
      xtick={-\r,-\a,0,\a,\r},
      % xticklabels={-\r,-\a,0,\a,\r}
      ytick={-\r,-\yint,0,\yint,\r}, 
      yticklabels={-\r,$-\sqrt\yintb$,0,$-\sqrt\yintb$,\r},
      x tick label style={anchor=north east}, % rótulos abaixo e a esquerda do eixo X
      y tick label style={anchor=south east}, % rótulos acima e a esquerda do eixo y
    ]

    % --- Curvas ---

    \draw[thick, blue] (axis cs:0,0) circle [radius=\r];    % Círculo de raio r

    \addplot[thick, red, domain=-1.2*\r:1.2*\r] ({\a}, {x});      % Reta vertical x=a
    \addplot[thick, red, domain=-1.2*\r:1.2*\r] ({-\a}, {x});     % Reta vertical x=-a

    % --- Região sombreada com Plot ---
    \path[fill=gray!30, opacity=0.6]
      plot[domain=(180-\ang):\ang, samples=100] ({\r*cos(\x)},{\r*sin(\x)}) --
      plot[domain=-\ang:-(180-\ang), samples=100] ({\r*cos(\x)},{\r*sin(\x)})
      -- cycle;          % Conecta o ponto inicial ao final


    % --- Linhas Tracejadas ---
     \addplot[Tracejado] coordinates {(-\a,{-\yint}) (\a,{-\yint})};
     \addplot[Tracejado] coordinates {(-\a,{\yint}) (\a,{\yint})};


    % --- Pontos de interseção ---
    \addplot[only marks, black!60, mark=*] coordinates {(-\a,{\yint})
                                                        (-\a,{-\yint})
                                                        (\a,{-\yint})
                                                        (\a,{\yint})
    };


  \end{axis}
  
\end{tikzpicture}
